\worldchapter{Götter und Glaube}{pantheon}
\label{ch:belief}

\begin{multicols*}{2}

\section{Götter und Glaube}\label{sec:gotter-und-glaube}

Tirakans Götterwelt ist vielseitig und omnipräsent.

Elf göttliche Geschwister wurden im ersten Zeitalter ihrer selbst gewahr.
Eine Blutfehde unter den Geschwistern sorge dafür, dass diese Götter sich auftrennten, und die Verräter entstanden.
In den folgenden Zeitaltern sollten Dämonen und Titanen auf der Welt erscheinen.
Sie wurden von Göttern und Verrätern geschaffen und sollten viele Jahrhunderte das Bild Tirakans prägen, bis sie aus unterschiedlichen Gründen in ihre Sphären getrieben wurden.

Zudem gibt es viele unabhängige Entitäten, Wesenheiten im Gefüge der Götter, welche frei existieren oder von anderen Kulturen als den Menschen verehrt werden.

Reisende von fernen Welten erreichten zu Zeiten des ersten Zeitalters Tirakan, um Wache über den Zwist der Götter zu halten.
Niemand, vermutlich nicht einmal sie selbst vermögen ihre Motivation zu beschreiben.
Später sollten Sie ein Volk nach ihrem Bild schaffen, welche gottgleiche Macht haben, die Drachen.

\section{Der Gläubige}\label{sec:der-glaubige}

Es steht jedem Wesen Tirakans frei, sich einer Gottheit zuzuwenden, dafür sind keine besonderen Voraussetzungen erforderlich.

Hat der Charakter mindestens einen Punkt in die Fertigkeit \textbf{Ritus} investiert, so darf er Anrufungen sprechen.

Jeder Charakter hat eine gewisse \textbf{Gunst}, die für die Beziehung zwischen ihm und seiner bevorzugten Gottheit steht.
Gunst wird für Anrufungen verwendet und kann durch gottgefälliges Handeln erlangt werden.

\subsection{Auslöser}
Belief wird genutzt, wenn ein Charakter eine Macht, Heilige, Ahnen oder namenlose Kräfte um Eingriff bittet.

\subsection{Petition-Ablauf}
\begin{enumerate}[leftmargin=*]
\item Bitte, Zeugin und Opfer benennen
\item \textbf{Gunst} ausgeben
\item optional +1 Burden für eine verzweifelte Bitte nehmen
\item \textbf{Petition}\index{Petition} werfen:
\[
 1d10 + Anrufungswert + Gunst + Modifikatoren
\]
\end{enumerate}

\[
 \text{Anrufungswert} = \text{Faith Level} + \text{Faith Rang}
\]

Maximale Gunst pro Anrufung:
\[
 \text{Maximale Gunst pro Anrufung} = 1 + \lfloor \text{Will} / 2 \rfloor
\]

\subsection{Zeremonielle Modifikatoren}
\begin{itemize}[leftmargin=*]
\item +1 passendes Opfer, Gelübde oder Preis
\item +1 geweihter Ort, Relikt oder anerkannte Zeugin
\item +1 Chor, Helfende oder Gemeinde
\item -1 hastig, improvisiert oder unterbrochen
\item -2 Bitte widerspricht klar Lehre, Eid oder Natur der angerufenen Macht
\item Gesamtkappung: höchstens +2, niedrigstens -2
\end{itemize}

\subsection{Ergebnisbander}
\begin{description}[style=nextline,leftmargin=0pt]
\item[\textbf{1--7 Silence}]\index{Silence} keine direkte Intervention, nur Zeichen oder Schweigen
\item[\textbf{8--11 Sign}]\index{Sign} symbolische Antwort, Warnung, Hinweis oder kurzer Schutz
\item[\textbf{12--15 Intercession}]\index{Intercession} spürbare Wunderwirkung im Rahmen der Szene
\item[\textbf{16+ Theophany}]\index{Theophany} großer Eingriff, fremd und verstörend
\end{description}

Gunst wird immer verbraucht. Bei \textbf{Silence} unter Tabubruch gibt es +1 Burden. Bei \textbf{Theophany} wählt der Guide einen heiligen Preis: Gelübde, Stigma oder Omen-Schuld.

\section{Petitionsformen}
\subsection{Shock Prayer}\index{Shock Prayer}
\begin{itemize}[leftmargin=*]
\item Ausloser: unmittelbare Gefahr
\item Dauer: 1 Action
\item Kosten: 1--2 Gunst
\item Wirkung: unmittelbare Standfestigkeit, Flucht, Warnung, Schutz oder Wucht
\end{itemize}
Bei \textbf{Sign}: +Faith Level auf einen sofortigen Test oder eine eingehende Consequence um eine Stufe senken.\\
Bei \textbf{Intercession}: zwei solche Boons.\\
Bei \textbf{Theophany}: die Szene kippt deutlich zugunsten der Gruppe, aber der Preis folgt.

\subsection{Blessing Rite}\index{Blessing Rite}
\begin{itemize}[leftmargin=*]
\item Dauer: 10 Minuten
\item Kosten: 2--4 Gunst
\item Ziel: Person, Objekt, Schwelle, Grab, Wasserstelle oder kleine gottliche Verunreinigung
\end{itemize}
Sign wirkt bis zum Szenenende, Intercession bis zur nächsten Full Rest, Theophany kann hartnäckige Flüche brechen oder heiligen Boden schaffen.

\subsection{Lesser Request}\index{Lesser Request}
\begin{itemize}[leftmargin=*]
\item Dauer: 10--15 Minuten
\item Kosten: 2--6 Gunst
\item Ziel: eine begrenzte Wunderbitte wie enthüllen, schützen, führen, reinigen, beruhigen oder plagen
\end{itemize}
Bei \textbf{Sign} gibt es ein klares Omen, eine Warnung, einen Weg oder kurzen Schutz.\\
Bei \textbf{Intercession} geschieht ein konkretes Wunder auf Szenenmaßstab.\\
Bei \textbf{Theophany} kippt die ganze Szene zugunsten der Bittenden, aber in verstörender oder unheimlicher Form.

\subsection{Great Invocation}\index{Great Invocation}
\begin{itemize}[leftmargin=*]
\item Dauer: 30+ Minuten
\item Kosten: 4--8 Gunst
\item Voraussetzung: mindestens eine Zeugin und ein bedeutendes Opfer, Relikt oder Eid
\end{itemize}
Dies ist die Form für öffentliche Rettung, Urteil, Massenprotektion oder heiliges Strafgericht.
Bei \textbf{Sign} kommen Furchtzeichen, Gewissheit oder ein zwingender Auftrag über die Szene.\\
Bei \textbf{Intercession} greift ein großes Wunder über Menge, Stadtteilrand, Schlachtfeldkante oder entweihten Ort.\\
Bei \textbf{Theophany} manifestiert sich eine szenenbestimmende Rettung, Verwandlung oder Verdammnis und fordert trotzdem einen heiligen Preis.

\section{Verbotene Praxis und Relikte}
\begin{itemize}[leftmargin=*]
\item \textbf{Relikt}\index{Relikt}: speichert 1--3 Arcana oder Gunst und überträgt sie durch Berührung oder Ritus
\item \textbf{Verbotene Methode}\index{Verbotene Methode}: senkt Zauberkosten um 1, mindestens auf 1, verursacht aber 1 Verderbnis
\end{itemize}

\section{Die Götter}\label{sec:die-gotter}

Einst erwachten elf Geschwister, welche in einer Zeit lange vor Menschen, Elfen und Zwergen die Welt bevölkerten.
Sie waren Töchter und Söhne unbekannter Eltern, und sie existierten in Einheit bis zu einem großen Verrat.
Sieben Verräter unter ihnen spalteten sich ab, und so blieben vier Götter auf der hellen Seite.

\subsection{Chronar - Gott der Zeit}\label{subsec:chronar-gott-der-zeit}

Chronar ist der Älteste der Geschwister, er wacht über den Fluss der Zeit.
Er wird oft als oberster Gott dargestellt und ihm werden die größten Tempel und Kathedralen errichtet.

Chronar wird oft als alter Mann mit einem Stundenglas dargestellt, obwohl es keine einheitliche Darstellung gibt.

Der höchste Feiertag, der mit Chronar in Verbindung gebracht wird, ist der letzte Tag des Jahres, das Fest des Jahreswechsels.

\subsection{Nadal - Wächterin der Magie}\label{subsec:nadal-wachterin-der-magie}

Nadal, die jüngste der Geschwister ist es, die die Magie auf die Welt brachte, und die sie auch bis heute hütet.
Alle Magie der Welt fließt durch ihre Hände und wird von ihr geleitet, zumindest ist es das, was die Priester und Anhänger von Nadal predigen.
Die Tempel der Nadal sind Orte des Studiums und der Lehre der Magie, und es wird ein höherer Wert auf Ruhe und Abgeschiedenheit als auf Prunk und große Gottesdienste gelegt.

\subsection{Algor - Behüter des Lichts}\label{subsec:algor-behuter-des-lichts}

Die Darstellungen von Algor als Gott treten sehr selten auf, häufig wird er als das Symbol der Sonne auf den Schriften und Tempeln dargestellt.
Selten sieht man ihn als jungen Mann mit einem kristallenen Stab in der Hand abgebildet, noch seltener wird er als alter Mann mit einem schlichten Holzstab skizziert.
Diese Darstellungen stammen vornehmlich aus der Zeit zwischen dem 2.
und 6. Jahrhundert.
Algor wird am ersten Tag des neuen Jahres gefeiert, mit dem Fest des neuen Jahres, das direkt an die Feierlichkeiten des Jahreswechsels anschließt.

\subsection{Herbarin - Der Gefallene}\label{subsec:herbarin-der-gefallene}

Herbarin war der zweitjüngste der elf Geschwister, die von unbekannten Eltern dazu auserkoren wurden, die Götter der Welt Tirakan zu sein.
Im dritten Zeitalter wurde er tödlich von dem Dämon Brahas getroffen, und sank sterbend und von Hass und Verbitterung erfüllt auf das Eis der Ebenen zu Kaal wo sich seine Spur für immer verlor.

\section{Die Verräter}\label{sec:die-verrater}

Sieben Geschwister spalteten sich im ersten Zeitalter von den Göttern ab und wandten sich gegen diese und die Wesen, die auf der Welt existierten.
Missgunst und Abneigung führte sie auf einen dunklen Pfad, und sie folgten fortan ihren eigenen, dunklen Interessen.

\subsection{Thzularn - Dunkle Göttin des Hasses}\label{subsec:thzularn-dunkle-gottin-des-hasses}

Thzularn gilt als Auslöserin des Streites zwischen den Göttern.
Sie führte die dunklen Götter an und schürte das Feuer des Hasses.
Sie gilt als unberechenbar und abgrundtief grausam, und ihre Anhänger opfern ihr oft mehr als nur ihr Leben.
Im Gegenzug bietet sie ihnen große Macht über die Menschen und noch mehr, wenn sie bereit sind, genug zu geben.
Ihr Bruder und Geliebter ist Dogan, der Meister der schwarzen Magie.

\subsection{Dogan - Dunkler Gott der Magie}\label{subsec:dogan-dunkler-gott-der-magie}

Dogan, Thzularns Bruder und Geliebter, brachte in der Zeit der Dämonenzüge eine grausame, dunkle Magie, die die Kräfte der Dämonen vervielfachte und das Gleichgewicht der Welt bedrohte.
Auch wenn seine Kräfte denen der Nadal unterlegen sind, ist seine Kraft immer noch present, und seine Anhänger wissen von ihr Gebrauch zu machen.

Besonders die höheren Künste der Magie verleiten dazu, die Kräfte des Dogan zu benutzen, hat sich ein Begabter jedoch einmal dazu entschlossen, Dogan zu dienen, wird dieser ihm wohl nie wieder die Freiheit geben.

\subsection{Seth'Nra - Dunkle Göttin, Wächterin der Pforte}\label{subsec:sethnra-dunkle-gottin-wachterin-der-pforte}

Seth'Nra ist es, die die Toten empfängt, falls sie sich im Angesicht des Rhodd'hrom für den falschen Weg entscheiden.
Die Toten werden für immer zu ihren Sklaven, und zu den Sklaven der Nekrologen und Begabten, die sich für ihren Weg entschieden haben.

In alten Schauermärchen ist oft die Rede von Seth'Nras Sklaventreiber der Seelen, Uhruk.
Erbarmungslos treibt er die ihm Untergebenen zur Arbeit an und holt angeblich unartige Kinder nachts zu sich in seine Reihen …

\subsection{Rhodd'hrom - Dunkler Gott der verfluchten Seelen}\label{subsec:rhodd'hrom-dunkler-gott-der-verfluchten-seelen}

Vor langer Zeit, als die Dämonen schon lange in die Flucht geschlagen waren und langsam Ruhe auf der Welt einkehrte, gelang es Rhodd'hrom, den Platz an der Seelenwaage einzunehmen, und so über das Schicksal der Toten zu bestimmen.
Diese Prüfung ist es, die jeder Mensch nach seinem Tod auf sich nehmen muss, besteht seine Seele nicht gegen die Stärke des dunklen Gottes, so ist er für immer ein Sklave von Seth'Nra.

\subsection{Ranorh - Dunkler Gott des verdammten Heeres}\label{subsec:ranorh-dunkler-gott-des-verdammten-heeres}

Zu der Zeit, als die Dämonen in die Unterwelt vertrieben wurden, wurde auch Ranorh's Einfluß auf die Welt geringer.
Der Gott, der als Anführer der Dämonen selbst die Gestalt eines Dämons anzunehmen pflegt, hat eine bedeutende Anhängerschaft unter den Dämonologen.
Diener des Ranorh zu sein verspricht treue und starke Sklaven, jedoch auch eine sichere Unendlichkeit in der Unterwelt.

\subsection{Ishaa'Nra - Dunkle Göttin der Verführung}\label{subsec:ishaa'nra-dunkle-gottin-der-verfuhrung}

Ishaa'Nra ist nicht so mächtig wie ihr Bruder Chronar, und ihre Versuche, das Zeitgefüge nachhaltig zu stören, scheinen bislang vergebens gewesen zu sein.
Jedoch bietet sie ihren menschlichen Verehrern die Macht, die Zeit zu beeinflussen, wenn auch nur in geringem Maße.

Ihre Verführungskünste sind von nichts übertroffen, auch gewährt sie ihren Anhängern ein besonderes Charisma.

Als Ishaa'Nras Bote tritt Nughlhuu häufig auf, er gilt als der zeitlose Übermittler, der Verkünder des widernatürlichen, vorzeitigen Todes.

\subsection{Dhas'Garyll - Der Geflohene}\label{subsec:dhasgaryll-der-geflohene}

Der Verbleib von Dhas'Garyll ist ein ungelöstes Rätsel, das die Gelehrten bis heute beschäftigt.
Der dunkle Gott, der einst im Zentrum des dämonischen Einflusses stand, zog sich vor der endgültigen Vertreibung der Dämonen aus den bekannten Ebenen zurück.
Sein gegenwärtiger Aufenthaltsort oder Zustand ist unbekannt, aber es wird angenommen, dass er weiterhin auf der Welt existiert.
Die Form, in der er sich aufhält, scheint ein weltlicher Körper zu sein, der es ihm ermöglicht, unerkannt unter den Sterblichen zu agieren.


\section{Die Titanen}\label{sec:die-titanen}

Nachdem Ranorh dreizehn grausame Dämonen schuf, welche die Welt bevölkerten, erschufen die Götter die Titanen, um die Welt Tirakan zu vor denen schützen.
Die Titanen sind greifbar für die Menschen und Wesen Tirakans, sie wandeln auf der Welt.
Oft wurden sie bei ihrem Wirken beobachtet.

\subsection{Ce-Nya - Titanin des Nebels}\label{subsec:ce-nya-titanin-des-nebels}

Ce-Nya ist wohl die geheimnisvollste Titanin.
Über sie gibt es viele Gerüchte, und vieles ist noch völlig unbekannt.
Mit ihr wird keine Gestalt verbunden, nur die Wölfe, ihre Diener, sind sehr häufig anzutreffen.
Ce-Nya gebietet über Tod, den Schlaf und das Vergessen und ist Schutztitanin der Diebe und Spitzbuben.
Ihre Tempel liegen versteckt und sind häufig für Fremde der Stadt nicht einfach zu erkennen.

\subsection{Dilae - Titanin der Magie}\label{subsec:dilae-titanin-der-magie}

Dilae ist nach Ce-Nya die geheimnisvollste der Titanen.
Sie dient Nadal im Sinne der Magie, und ihre Tempel sind, ähnlich denen der Nadal, eher mit Bibliotheken denn mit mit zeremoniellen Orten vergleichbar.
Ihre Dienerwesen sind wohl die seltsamsten Wesen, die auf der Welt vorkommen.
Die Bolde gibt es in allen nur erdenklichen Gestalten, und sie sind in ihrer Art unberechenbar.
Dilae gebietet über die Magie und ist Schutztitanin der Begabten.

\subsection{Ginae - Titanin des Wassers}\label{subsec:ginae-titanin-des-wassers}

Ginae, die Titanin des Wassers, wurde von Chronar geschaffen, und schloss die Morgalas für immer in den Tiefen des Schlundes nahe der Ebene von Jaar ein.
Ginae gilt als aüßerst launisch und wechselhaft in ihren Handlungen auf der Welt.
Sie ist die Schutztitanin der Seeleute und gebietet über Geburt, Neuentstehung und Erneuerung.

Ihre Dienerwesen sind die Hydren, jene mehrköpfigen Wesen welche es in allen Meeren Tirakans in den Tiefen gibt.
Allen voran steht Hydonia, die legendenbehaftete große Hydra aus dem Felsenmeer.

\subsection{Gronar - Titan der Erde}\label{subsec:gronar-titan-der-erde}

Auch von Gronar ist wenig bekannt.
Gronar dient dem Humus und der Erde, und gilt als ein sanftmütiger, den Menschen zugetaner Titan.
Seine Diener sind die Riesen, und auch wenn die Menschen sie sehr selten zu Gesicht bekommen, gibt es wohl mehr, als man annimmt.
Gronar gebietet über das Leben selbst und ist Schutztitan der Bauern.

\subsection{Jogran - Titanin des Eises}\label{subsec:jogran-titanin-des-eises}

Jogran steht für das Eis und den Winter.
Jogran wird nachgesagt, sie hätte einen Eispalast weit im Norden, in einem Land, das von grausamen großen Wesen bewohnt wird.
Dort sei auch der Ursprung ihrer Diener, der Einhörner.
Sie gebietet über den Winter und das Eis und ist Schutztitanin der Jäger.

\subsection{Nilan - Titan des Lichts}\label{subsec:nilan-titan-des-lichts}

Nilan ist der Titan des Lichts, und die Menschen bauen ihm große, helle Tempel, die denen des Algor ähnlich sind.
Ihm dienen die geflügelten Schlangen, die weit draußen auf der See heimisch sind und nur selten über dem Land zu sehen sind.
Nilan wird dem Recht und dem Gesetz zugesprochen und ist Schutztitan aller Paladine.

\subsection{Pailos - Titan der Luft}\label{subsec:pailos-titan-der-luft}

Pailos, Herr über die Luft, das Wetter und die Winde, gilt als ein gütiges Wesen, dass im Dienste der Götter die Winde der Welt lenkt und anfacht.
Er gebietet über das Wetter und die Wetteränderungen und ist Schutztitan der Gaukler, der Hochstapler und der Scharlatane.

Seine Diener sind die Pegasi, geflügelte Pferde, die in den großen Steppen des Nordens anzutreffen sind.

\subsection{Rogal - Titan des Erzes}\label{subsec:rogal-titan-des-erzes}

Rogal ist der Titan des Erzes, Schutztitan der Schmiede.
Seinem Sohn Taurus dienen die Minotauren, menschenähnliche Wesen mit Köpfen gleich denen der Stiere.
Die Minotauren leben tief unter der Oberfläche, und man sagt ihnen nach, dass sie die einzigen Wesen seien, die Kontakt zu den verfluchten Morgalas haben.

Rogal wird mit Standhaftigkeit und Treue in Verbindung gebracht.

\subsection{Tador - Titan des Steins}\label{subsec:tador-titan-des-steins}

Der mächtige Tador, Schutztitan der Bergleute, ist der Vater des Steins.
Er gebietet über die Sicherheit und den Schutz.
Die Dienerwesen des Tador sind die mächtigen Steinwesen Tonarr und Astaron, der im Jahr 722 EC. im Kampf gegen Telatoon beim wahren Sturm des Hufen fällt.
Ebenfalls dienen ihm die irdischen Centauren, große Wesen mit Körpern gleich denen der Pferde und Oberkörpern und Köpfen von Menschen.
Die Centauren leben, wie auch die Minotauren, größtenteils unter der Erde.
Nur selten wagen sie sich auf die Oberfläche.
Die Eingänge zu ihren Höhlen liegen versteckt in den Wäldern des Nordens.

\subsection{Zyral - Titanin des Feuers}\label{subsec:zyral-titanin-des-feuers}

Zyral brachte den Wesen auf der Welt das Feuer, und rüstete sie zum Kampf.
Geschaffen von Chronar ist Zyral die Schutztitanin der Krieger und kämpfenden Klassen, und sie gebietet über Waffen und den erfolgreichen Kampf.

Die Greifen dienen Zyral in ihrem Auftrag, jene großen gefiederten Wesen, deren Leib der eines Löwen, deren Kopf jedoch der eines Adlers ist.
Sie folgen Fessaal, dem Feuergreifen aus dem Kemalgebirge.
Ihre Gefährtin ist Phidia, die Feuergreifin.
Ihr zu Ehren hat sich auch der Orden der Phidia gegründet deren Motto ``Feuer, Flamme Schwert'' ist.

\section{Erzdämonen}\label{sec:erzdamonen2}

Im zweiten Zeitalter von Ranorh geschaffen, haben die Dämonen lange die Welt für sich beansprucht.
Erst durch das Wirken Dhas'Garylls sind sie im dritten Zeitalter in die Dämonenebene vertrieben worden.
Diese titanengleichen Wesen bevölkern nun die Unterwelt.

Erzdämonen haben direkte Dienerwesen, die üblicherweise in ihrer Nähe zugegen sind, aber auch von ihnen auf die Welt entsandt werden.

\subsection{Duglaraan - Der verfluchte Ersäufer - I}\label{subsec:duglaraan-der-verfluchte-ersaufer}

Die erste Domäne ist die Welt des Wassers, über die Duglaraan mit äußerster Grausamkeit herrscht.
Duglaraan, der verfluchte Ersäufer und Schänder der Seefahrt, hat die Gestalt einer riesigen Seeschlange mit sieben Köpfen und großen Hörnern auf der Rückseite der Wirbel.
Seine beiden Dienerwesen, Shuugai und Ganod, sind nicht besser als er, sie sind seine Begleiter und haben eine sehr ähnliche Gestalt, wenn auch wesentlich kleiner.

\subsection{Calysporn - Der verfluchte Herrscher des Nebels - II}\label{subsec:calysporn-der-verfluchte-herrscher-des-nebels}

Über diese Domäne herrscht Calysporn, der verfluchte Herrscher des Nebels und der Nacht.
Calysporn, der keine feste Gestalt hat, wird begleitet von Dugai und Sarag'Roh, die beide eine menschenähnliche Gestalt haben, wenn auch ihre Körper nur aus Nebel bestehen.

Dugai scheint männliches Aussehen zu haben, während Sarag'Roh wohl seine Gefährtin ist.
Sie vermögen es des nächtens unter Türen und durch Fenster hindurch zu gleiten und werden im Volksmund verdächtigt, Kinder in ihren Krippen ums Leben zu bringen.
Dem Dämonologen können sie die Gabe der Unsichtbarkeit oder des Nebelkörpers verleihen.

\subsection{Gyral - Der verfluchte Entzünder - III}\label{subsec:gyral-der-verfluchte-entzunder}

Die dritte Domäne ist die Welt des Feuers.
Über diese Welt herrscht Gyral, auch Gyralph genannt, mit seinen furchtbaren Dienern.
Gyralph, eine große Gestalt gleich einem Troll, trägt den Kopf eines Adlers.
An der Stelle des Schnabels befinden sich drei lange Rüssel, die bis zum Boden reichen.
Gyralph scheint immer zu brennen, seine Haut gleicht glühender, schwarzer Asche mit unzähligen feurigen Wunden.
Seine Diener, Nugshu und Tulhu, gleichen den Greifen, jedoch tragen sie vier Hörner auf dem Rücken und die gleiche brennende Haut wie ihr Meister.

\subsection{Cthupheen - Der Tod aus der Erde - IV}\label{subsec:cthupheen-der-tod-aus-der-erde}

Wer immer Cthupheen begegnet, wird diese Begegnung sein Leben lang nicht vergessen.
Viele hundert Schritt misst der gigantische, irdene Wurm, an dessen Maul tausende von Zähnen gleich denen des Haifisches prangen.
Wenn auch seine Begleiter, Ijothu und Shoggu Hai, um vieles kleiner sind, so haben sie auch noch eine Länge von mehreren Schritt und einen Durchmesser von gut einem halben Schritt.
Die schwarzen Würmer gelten als der Tod aus der Erde und verschlingen oder erwürgen ihre Opfer.

\subsection{Brahas - Vater der Sucht - V}\label{subsec:brahas-vater-der-sucht}

Über die fünfte Domäne herrscht Brahas, der den sterbenden Gott Herbarin gesteinigt hat.
Er wird als der Vater der Sucht und der Schänder der Freuden bezeichnet.
Brahas Diener sind Yoghas und Bhulu, zwei steinernde Wesen in der Gestalt von Trollen mit toten, schwarzen Augen.
In seiner Umgebung werden alle lebenden Wesen träge, faul und schläfrig.
Kaum ein Wesen, außer Priestern und geweihten oder heiligen Tieren, kann sich dem entziehen.
Wie Rauschmohndampf legt sich der Atem Brahas um die, welche in seiner Nähe sind.
Yoghas und Bhulu allerdings sind trotz ihrer steinernen Gestalt und anscheinender Abwesenheit von Leben sehr kräftige Wesenheiten.
Sie vermögen es jeden den sie berühren in einen ewigen, fast schlafartigen Rausch zu versetzen indem das Opfer durch Alpträume solange gequält wird bis nach einigen Aeonen der Tod eintritt.
Allerdings hat dies noch kein Mensch dokumentieren können, sodass diese Qual auch ewig dauern kann.

\subsection{Lhugghu - Der verfluchte Herr der Katastrophen - VI}\label{subsec:lhugghu-der-verfluchte-herr-der-katastrophen}

Lhugghu, der diese Domäne anführt, ist der dämonische Herr des Windes und des Sturms, der Erschaffer des Chaos und des Neides.
Lhugghu hat die Gestalt einer schwarzen Wolke ohne ein Gesicht, und seine Diener, Zzibra und Cayogbaar, begleiten ihn in Form von mißgestalteten Raben ohne Flügel und mit schwarzer Haut.

\subsection{Haiphen - Der verfluchte Bringer des Todes mit dem Stahl - VII}\label{subsec:haiphen-der-verfluchte-bringer-des-todes-mit-dem-stahl}

Die siebte Domäne wird beherrscht von Haiphen, der die menschliche Form eines Kriegers in schwarzer Rüstung, jedoch ohne ein Gesicht hat.
Haiphen, der als verfluchter Herr des Erzes und Bringer des Todes mit dem Stahl bezeichnet wird, trägt ein blau glühendes Schwert und einen brennenden Streitkolben, mit dem er tödliche Wunden schlägt.
Seine Begleiter sind Ghyogh und Hasoo, die schwarzen Bluthunde ohne Augen.
Er gilt vielen Söldnern, welche sogar dem Echsenheer oder den Minotauren dienen, als Schutzherr.

\subsection{Lothcam - Der verfluchte Beherrscher der Schatten - VIII}\label{subsec:lothcam-der-verfluchte-beherrscher-der-schatten}

Lothcam, der Beherrscher der Schatten und Schattenwesen hat die Gestalt einer riesigen Amöbe oder eines riesigen Wurms.
Seine Diener sind Thurgroth und Brasho.
Lothcam folgt seinem eigenen Plan.
Er will nichts anderes als die Herrschaft über ganz Tirakan.

\subsection{Nughlhuu - Der Bote Ishaa'Nras - IX}\label{subsec:nughlhuu-der-bote-ishaa'nras}

Nughlhuu ist ein Vierbeiner, der einem Einhorn gleicht.
Schwarze Flüssigkeit trieft dauerhaft von seinem mageren, fast verhungerten Körper.
Man sagt, das der, der in das Licht blickt, dass von dem Horn auf seinem Kopf ausgeht, für immer zu Stein erstarrt und dämonische Qualen leiden muss, während er von Nycalhu und Nugaa langsam zerfressen wird.
Die beiden Begleiter Nughlhuus haben die Gestalt skelettierter Vögel, von deren bleichen Knochen die modernden Federn herabhängen.

\subsection{Iiinyca - Der Schänder der magischen Kraft - X}\label{subsec:iiinyca-der-schander-der-magischen-kraft}

Iiinyca, der Schänder der Magischen Kraft, ist der Herrscher über die zehnte Domäne, die dämonische Welt der Magie.
Iiinyca tritt selbst nicht in Erscheinung und ist nur als unglaublich starke dämonische Magie zu erahnen, die alles magische in ihrer Umgebung stört und selbst von nicht Begabten wahrgenommen wird.
Iiinycas Begleiter, Tur Oggu und Yogzzi, sind dunkle, grausame Bolde, die mit ihren scharfen Zähnen und mit unheiliger Magie starke Gegner für jeden Rechtschaffenden darstellen.

\subsection{Nycalha - Die verfluchte Herrin der perversen Lust - XI}\label{subsec:nycalha-die-verfluchte-herrin-der-perversen-lust}

Die elfte Domäne wird angeführt von Nycalha, die die dämonische Eifersucht und die perverse Lust bringt.
Nycalha's Gestalt ist die einer großen Frau mit bleicher, kranker Haut, die kein Haupthaar hat.
Ihre Begleiter, Bradhi und Flamer, treten als Zwillingspaar auf, beides wunderschöne Frauen, die eine mit flammendem roten Haar und bleicher Haut, die andere mit schwarzem vollen und sinnlichen Haar mit südländischer Haut.

\subsection{Nyphen - Der Herr des kalten Todes - XII}\label{subsec:nyphen-der-herr-des-kalten-todes}

Nyphen, der verfluchte Herr des Eises und des kalten Todes, ist eine kristallene, tiefblaue Gestalt mit tausenden von Dornen und Stacheln.
Die riesenhafte, eiskalte Wesenheit wird begleitet von Lhuu und Baarhai, schwarzen, krötenförmigen Dämonen.
Deren Stacheln, vom Körper geschossen, reissen furchtbare Wunden die niemals heilen.

\subsection{Wisgu - Gebieter des Ersten Bundes - XIII}\label{subsec:wisgu-gebieter-des-ersten-bundes}

Wisgu, Hasser der Elfen und Zwerge, Gebieter des ersten Bundes ist eine nicht greifbare Gestalt aus Staub und reiner Energie.
Im dritten Zeitalter ging Wisgu den Bund mit den Morgalas ein, und versklavte sie zu immerwährender Dienerschaft.
Seine Diener Turzzi und Ylvath haben die Gestalt irdener Statuen.


\section{Die Reisenden}\label{sec:die-reisenden}

Die Götter waren sich ihrer selbst noch nicht bewusst geworden, als vier dieser urzeitlichen Wesen auf der Welt erschienen.
Niemand wusste, woher sie kamen.
Den Grund ihrer Existenz konnten nur sie selbst bestimmen.
Die Götter bemerkten sie erst spät, als sie die Welt schon eine Weile beobachtet hatten.

Es waren vier Reisende, die aus der Ferne kamen.
Ihre Macht war der der Götter gleich, und sie sollten die Geschichte Tirakans in den kommenden Zeitaltern entscheidend mitbestimmen.

\subsection{Hagarun - Die die Liebe und den Hass eines ganzen Volkes in sich trägt}\label{subsec:hagarun-die-die-liebe-und-den-hass-eines-ganzen-volkes-in-sich-tragt}

Hagarun ist die Reisende der Emotionen, ein Wesen, das die tiefsten Gefühle der Sterblichen in sich aufnimmt.
Ihre Erscheinung ist wandelbar und spiegelt die Seelen derer wider, die sie betrachten.
Mal erscheint sie als strahlende, mal als schattenhafte Gestalt, je nachdem, welche Emotionen sie in sich trägt.

Ihr Einfluss ist subtil, aber allgegenwärtig.
Sie knüpft Bande zwischen Menschen, stärkt Bündnisse und entfacht Feindschaften.
Viele Kriege und große Liebesgeschichten Tirakans werden ihr zugeschrieben, als hätte sie die Herzen der Beteiligten gelenkt.
Ihre Macht liegt nicht in der direkten Manipulation, sondern in der Verstärkung dessen, was bereits in den Sterblichen schlummert.

\subsection{Irminar - Der Ruhe ist}\label{subsec:irminar-der-ruhe-ist}

Irminar verkörpert Stille und Unveränderlichkeit als Gegenpol zu Modordes ständiger Bewegung.
Er ist der Reisende des Gleichgewichts, dessen Gegenwart als unendliche Nacht empfunden wird.
In den wenigen Überlieferungen, die ihn beschreiben, wird er als ein schweigendes, in Nebel gehülltes Wesen dargestellt, dessen Augen die Ewigkeit zu durchdringen scheinen.

Seine Macht liegt im Stabilisieren und Bewahren.
Während Modorde Veränderungen einleitet, sorgt Irminar dafür, dass das Fundament intakt bleibt.
Sein Wirken zeigt sich vor allem in Momenten, in denen das Chaos die Welt aus den Fugen zu reißen droht.
Viele sehen in ihm die unsichtbare Hand, die Tirakan im Gleichgewicht hält.

\subsection{Modorde - Der das Sein beendet und neu erschafft}\label{subsec:modorde-der-das-sein-beendet-und-neu-erschafft}

Modorde ist ein Reisender des Wandels, der Zerstörung und der Wiedergeburt.
Sein Wesen ist geprägt von einem tiefen Verständnis des Vergehens und des Neubeginns.
Oft erscheint er als schattenhafte Gestalt, deren Form sich ständig wandelt - ein Spiegelbild seiner Natur.
Seine Anwesenheit wird von einem Gefühl des Endes begleitet, das Furcht und Hoffnung zugleich weckt.

Er ist kein grausamer Zerstörer, sondern ein unerbittlicher Erneuerer.
Für Modorde ist jedes Ende der Beginn eines neuen Zyklus.
Sein Einfluss ist besonders in Zeiten großer Umbrüche spürbar, wenn alte Strukturen zusammenbrechen und Platz für Neues geschaffen wird.
Die Legende erzählt, dass Modorde selbst das Zeitalter der Drachen beendete, um Platz für die Völker zu schaffen.


\subsection{Odovacar - Der Feuer und Leben aus der Ferne bringt}\label{subsec:odovacar-der-feuer-und-leben-aus-der-ferne-bringt}

Odovacar ist der Reisende des Feuers, ein ungezähmter Geist, der sowohl Leben als auch Zerstörung bringt.
Er wird oft als brennende Gestalt beschrieben, ein wandelnder Sturm aus Flammen und Hitze.
Wo er wandelt, entstehen sowohl fruchtbare Böden als auch verwüstete Landstriche - ein Sinnbild für seine zwiespältige Natur.


\section{Die Drachen}\label{sec:die-drachen}

Geboren als Kinder der Reisenden sind die Drachen die wohl mächtigsten unabhängigen Wesen Tirakans.
Die Reisenden kamen zu einer Zeit auf die Welt, da die Götter noch nicht einmal ihres Ursprungs gewahr waren, und beobachteten die Geschehnisse, ohne in die Kämpfe der Vorzeit einzugreifen.
Zu Ende des ersten Zeitalters begab es sich, dass die Reisenden sieben Kinder zeugten, um die Kämpfe der Titanen und Dämonen zu überwachen.
Diese Kinder waren fortan Drachen genannt.
Den Legenden nach leben vier dieser sieben Kinder noch heute auf Tirakan, wenngleich auch nur wenige, nicht bestätigte Berichte von ihnen erzählen.

\subsection{Tar - Der Mutige}\label{subsec:tar-der-mutige}

Tar der Mutige war das erste Kind der Reisenden und ist somit der Älteste der sieben Geschwister.
Tar, der auch der mächtigste der Drachen ist, hält sich bis heute hauptsächlich in den nördlichen Bergen auf.
Von den Elfenvölkern als Halbgott verehrt, nimmt er bei den Atiarel den Platz des höchsten Gottes ein.
Nach dem Drachen Tar ist auch ein Fluss benannt, der von den Bergen der Atiarel durch die Ebenen von Meridian bis zum Felsenmeer fließt.

\subsection{Aspersia - Die Gefallene}\label{subsec:aspersia-die-gefallene}

Während des Kampfes am Schlund wurde Aspersia mit den Dämonen hinabgezogen.
Ihr Schicksal ist nicht bekannt.
Legenden zufolge soll sich Aspersia jedoch in der Erschaffung eines eigenen Volkes versucht haben.
Aufgrund der tragischen Umstände und der dämonischen Umgebung muss es dabei jedoch zu einem furchtbaren Fehler gekommen sein.
So erschuf Aspersia zwei Völker nach ihrem Bilde.
Die eine in humanoider Gestalt, die andere ihrer Form gleich.
Diese wurden fortan Sethlarn genannt.

\subsection{Derumir - Der Rote}\label{subsec:derumir---der-rote}

Derumir, „der Rote“, ist einer der sieben Urdrachen, die von den Reisenden geschaffen wurden, um über die Welt zu wachen.
Seine leuchtend roten Schuppen und die feurige Aura, die ihn umgibt, machen ihn zu einem Symbol elementarer Urkraft.

Er zeigt sich nur in Zeiten äußerster Not, wenn die Grundfesten Tirakans bedroht sind, und sein Erscheinen ist selten ein Segen für die Sterblichen.
Wo er wandelt, folgt ihm das Feuer als Zeichen seiner untrennbaren Verbindung mit den vulkanischen und chaotischen Kräften der Welt.
Sein Einfluss ist meist indirekt, aber in den Gebieten, in denen er vorkommt, spürbar.
Er bevorzugt die Nähe von brodelnden Vulkanen und urzeitlichen Lavaströmen, und seine Anwesenheit kann ganze Landstriche verändern.

\subsection{Saranor - Der Schnelle}\label{subsec:saranor-der-schnelle}

Saranor der Schnelle zog sich aufgrund einer Auseinandersetzung mit den Reisenden von der Oberfläche Tirakans zurück.

\subsection{Tedasiél}\label{subsec:tedasiel}

Tesasiél geriet während der Zeit der Dienerwesen in einen Streit mit Irminar und wurde von ihm getötet.
Die Legenden erzählen davon, dass sein Tod mit dem Auflehnen der Drachen gegen die Reisenden zu tun hat.

\subsection{Samusa - Die Mutter}\label{subsec:samusa-die-mutter}

Die älteste Tochter der Hagarun ist die Mutter des Volkes der Elfen.
Als die Drachen den Verrat an ihren Eltern begingen war sie es, die die Elfen gebar.
Seither ist, unter Anderem in der Iana Alyaria, von Samusas Hort die Rede, wenn es um den Ritus geht.

\subsection{Héroda - Der aus der Nacht kommt}\label{subsec:heroda-der-aus-der-nacht-kommt}

Héroda ist einer der sieben Urdrachen von Tirakan, die von den Reisenden am Ende des ersten Zeitalters geschaffen wurden.
Wie seine Geschwister ist Héroda ein göttliches, unsterbliches Wesen, das sich von der Welt der Sterblichen fernhält.
Er zeigt sich nur in seltenen Momenten, wenn das Gleichgewicht der Schöpfung selbst bedroht ist.
Anders als Tar der Mutige oder Aspersia, deren Sturz in den Schlund legendär ist, vermeidet Héroda jede Form der direkten Einflussnahme.
Seine Existenz ist eine Konstante im Verborgenen, eine stille Wache über die Schattenseiten der Welt.


\section{Unabhängige Entitäten}\label{sec:unabhangige-entitaten}

Unabhängige Entitäten göttlicher Natur gibt es viele auf Tirakan.
Völker wie die Quitaron oder Xordai beten ihre eigenen, den Reichen fremden Gottheiten an, Wesenheiten wie Robärthrunius haben selbst einen gottähnlichen Status, wenngleich auch ohne Anhänger.


\subsection{Argorin - Der Erschaffer}\label{subsec:argorin-der-erschaffer}

Die Xordai unterscheiden sich in ihrem Glauben von den meisten Völkern Tirakans.
Diese verehren in vielfältigen Formen die vier bzw.
sieben Götter, die auch von den Menschen als höchste Wesen der Welt angesehen werden.
Die Xordai jedoch beten zu den beiden Göttern Argorin und Dargumir.

Argorin ist der ehrfürchtige, der das Leben in all seinen Formen achtet.
Er ist der Erschaffer, der die Glut im Innern der Welt am Leben hält.
Seine Taten sind weise, und er kennt seinen Weg hin zu einer besseren Welt ohne die fürchterlichen Sethlarn.


\subsection{Dargumir - Der Narr}\label{subsec:dargumir-der-narr}

Dargumir ist der Narr.
Er ist Argorins Gefährte und begleitet ihn auf der immerwährenden Suche.
Anders als Argorin ist er fern von einem inneren Frieden.
Sein Handeln ist Chaotisch, seine Werke zum Teil grausam.
Er ist die Hoffnungslosigkeit, steht aber auch für das Leben im absoluten Jetzt.
Dargumir wird die Zerstörung und der Tod zugeschrieben.

\subsection{Akhrosch - Der Gott der Orks}\label{subsec:akhrosch-der-gott-der-orks}

Der erste Vater war Akhrosch.
Damals, kurz nachdem Licht von Dunkelheit durch den Allvater und die Allmutter getrennt wurde, erschien er auf Tirakan.
Er wurde der Ungeborene genannt denn er hatte keinen Vater und keine Mutter.
Aus sich selbst und dem Lehm der Erde schuf er Kchtuh, seine Gefährtin.
Gemeinsam jagten sie die damaligen Lebewesen Tirakans und gebaren viele Kinder.

… von Edis Elbenfreund, Anthropologe zu Yavon, ca.300 EC

\subsection{Taxaros - Herr des Feuers}\label{subsec:taxaros-herr-des-feuers}

 Taxaros ist ein mächtiger Feuerelementarer, der, schätzungsweise im 8. Jahrhundert, durch Meister Astrum Fabula zur Akademie von Jaar nach Tirakan geholt wurde.
 Unklar ist allerdings, wie Fabula dies vollbringen konnte, da dieser Meister der Schwarzen Magie zu Jaar war.
 Man vermutet, dass er, bei dem mehrere Dekaden dauernden Ritual das nötig war, um Taxaros zu binden, Hilfe durch einen Elementaristen hatte.

\subsection{Robärthrunius - Mystischer Klabauter}\label{subsec:robarthrunius-mystischer-klabauter}

Robärthrunius ist ein Mythischer König der Klabauter, der auf einem verzauberten Schiff die Meere kreuzen soll.
Er wird von vielen als gottgleich bezeichnet, doch es gibt immer wieder Geschichten, die seine Sterblichkeit andeuten.
Einige Berichte erzählen von Wunden, die selbst er nicht sofort heilen konnte, oder von Momenten der Schwäche, die ihn kurz innehalten ließen.

Niemand weiß genau, warum Robärthrunius die Meere durchstreift, ohne Rast und ohne Ziel.
Manche sagen, er sei auf der Suche nach etwas Verlorenem, andere glauben, er sei auf der Flucht vor einer großen Gefahr.
Vielleicht teilt er mit den Stürmen mehr als nur die Liebe zum Meer - vielleicht ist er selbst ein Teil von ihnen.

\subsection{Hirrukh - Der, der das Rad der Schöpfung in Händen hält}\label{subsec:hirrukh-der-der-das-rad-der-schopfung-in-handen-halt}

Die Goblins der Asch-Ta-Khi verehren das Gottwesen Hirrukh, welches in ihrem Augen das Rad der Schöpfung in den Händen hält.
Hirrukh wird bei den verschiedenen Sippen unterschiedlich verehrt, mal als Heiler und Lebensspender, mal als wilder Krieger oder rachsüchtiger Vernichter.

\subsection{Caltae - Minotaurenschamanin}\label{subsec:caltae-minotaurenschamanin}

Caltae ist die Gefährtin von Taurus.
Sie war es, bevor die Minotauren für tausend Jahre in eine Lethargie verfielen, und sie wird es für immer sein.
Caltae lag mit Taurus in der Kammer des Schlafs, paralysiert und unfähig, sich zu regen oder Gefühle zu empfinden.
Mit dem Erwachen von Taurus im Jahre 0 EC erwachte auch sie, und sie sollte für die nächsten 1000 Jahre die gewissenhafte Seite von Taurus sein.

\subsection{Der Chronist}\label{subsec:der-chronist}

Der Chronist ist eine mythische Gestalt Tirakans, die angeblich die Geschichte der Welt in seiner Ganzheit aufzeichnet.
Einige sprechen von ihm als eine Wesenheit fern dieser Welt, andere behaupten es sei Chronar selbst der die Geschichte in seinem ewigen Buch festhält.
Neben dem Chronisten existiert auch noch die philosophische Vorstellung eines Erzählers.

Der Chronist ist dem Volk Tirakans unbekannt.
Die Gestalt des Chronisten und des Erzählers existiert allenfalls in magietheoretischen oder religiösen Werken großer Philosophen.
Sie spielen keinerlei Rolle im Geschehen der ersten 950 Jahre des fünften Zeitalters.

\subsection{Der Träumer}\label{subsec:der-traumer}

Es gelang ihm kaum, seine Freude zu verbergen.
Lange hatte er auf diesen Moment warten müssen.
All die unzähligen Jahre hatte der Chronist ihn für seinen Sohn gehalten, doch er wußte es besser.
Er war nicht zum Chronisten geboren.
Nein, er war ein Träumer.
Während sein "Vater" die Geschichte Tirakans aufzeichnete, wie sie passierte, sah er ihm über die Schulter und träumte davon, wie sie wohl weitergehen mochte.
Wie erstaunt er doch war, als er feststellte, daß von Zeit zu Zeit eine seiner hämischen, kleinen Ideen ihren Weg aus seinen Träumen in die Welt Tirakans fand.
Dies geschah recht selten, ein Frevel hier, ein dämonischer Pakt dort - Kleinigkeiten in seinen Augen.
Doch war es für ihn ein Zeichen seiner Bestimmung.

\end{multicols*}
